% --- Archon physics formalization source begin ---
% archon:physics
% archon:covers PhyXMiniProblems/problem_phyx_mini_0757.lean
% archon:source-report reports/phyx_mini/problem_phyx_mini_0757.source.json

\chapter{Physics problem phyx\_mini\_0757}
\label{ch:PhyXMiniProblems_problem_phyx_mini_0757}

\paragraph{Problem source.}
\#\# Physical scenario

Figure shows an arrangement in which four disks are suspended by cords. The longer, top cord loops over a frictionless pulley and pulls with a force of magnitude \$98\textbackslash{} N\$ on the wall to which it is attached. The tensions in the three shorter cords are \$T\_1 = 58.8\textbackslash{} N\$, \$T\_2 = 49.0\textbackslash{} N\$, and \$T\_3 = 9.8\textbackslash{} N\$.

\#\# Question

What is the mass of disk \$C\$?

\#\# Answer choices

- A: A: 1.0 kg

- B: B: 2.0 kg

- C: C: 4.0 kg

- D: D: 5.0 kg

\#\# Figure caption (auxiliary; use the image as primary evidence)

The image depicts a mechanical system involving a pulley and a series of four green spheres labeled A, B, C, and D, connected vertically by strings. 

- The system is anchored to a structure on the left with a pulley at the top right.
- The string runs over the pulley and supports the four spheres.
- The tension in the strings between the spheres is indicated as \textbackslash{}(T\_1\textbackslash{}), \textbackslash{}(T\_2\textbackslash{}), and \textbackslash{}(T\_3\textbackslash{}), in descending order.
- Each sphere is connected by the string, suggesting a downward force due to gravity, and the system demonstrates a pulley problem often used in physics.
  
The structure is colored brown, and the pulley is depicted in blue.

\#\# Dataset metadata

Dynamics; Physical Model Grounding Reasoning, Multi-Formula Reasoning

\paragraph{Recorded answer/context.}
C: C: 4.0 kg

\paragraph{Figure/image path.}
/root/proposal\_for\_physic/science-mango/phyx\_mini\_run/phyx\_data/test\_image/757.png

\paragraph{Formalization target.}
create a compiling Lean file with sorry bodies at `PhyXMiniProblems/problem\_phyx\_mini\_0757.lean`. The Lean declarations must preserve the physical quantities, dimensions or dimensional roles, figure labels, governing-law hypotheses, and final relation expressed by this problem.
Use Mathlib/Physlib names found through LeanExplore where available. If a physics API is missing, introduce faithful local abstractions rather than scalar placeholder aliases.

\begin{theorem}[Physics formalization target]
\label{thm:physics:phyx_mini_0757:target}
\lean{PhyXMiniProblems.ProblemPhyXMini0757.problem_phyx_mini_0757}
\uses{def:physics:phyx-mini-0757:phyxminiproblems-problemphyxmini0757-massinkilograms, def:physics:phyx-mini-0757:phyxminiproblems-problemphyxmini0757-fourdiskpulleysetup, def:physics:phyx-mini-0757:phyxminiproblems-problemphyxmini0757-matcheswrittenscenario, def:physics:phyx-mini-0757:phyxminiproblems-problemphyxmini0757-matchessuppliedfigure, def:physics:phyx-mini-0757:phyxminiproblems-problemphyxmini0757-matchesproblemforcereadouts, def:physics:phyx-mini-0757:phyxminiproblems-problemphyxmini0757-usesstandardnearearthgravity, def:physics:phyx-mini-0757:phyxminiproblems-problemphyxmini0757-hasphysicalparameters, def:physics:phyx-mini-0757:phyxminiproblems-problemphyxmini0757-satisfiessuspendeddiskstaticslaws, def:physics:phyx-mini-0757:phyxminiproblems-problemphyxmini0757-isuniquematchingdisplayedmass}
The assigned autoformalize agent should translate this physics problem into Lean declarations in the covered file, with theorem and lemma proof bodies written as `by sorry`.
\end{theorem}
\begin{proof}
This is an autoformalization task, not a proof task. Produce statements that can later be proved by the physics prover without weakening the physical model.
\end{proof}

% --- Archon declaration topology begin ---
\section{Declaration topology}
\begin{definition}[acceleration Dimension]
  \label{def:physics:phyx-mini-0757:phyxminiproblems-problemphyxmini0757-accelerationdimension}
  \lean{PhyXMiniProblems.ProblemPhyXMini0757.accelerationDimension}
  The physical dimension L T⁻² of an acceleration magnitude
\end{definition}
\begin{proof}
  This is the declared model definition of acceleration Dimension.
\end{proof}
\begin{definition}[force Dimension]
  \label{def:physics:phyx-mini-0757:phyxminiproblems-problemphyxmini0757-forcedimension}
  \lean{PhyXMiniProblems.ProblemPhyXMini0757.forceDimension}
  The physical dimension M L T⁻² of a force magnitude
\end{definition}
\begin{proof}
  This is the declared model definition of force Dimension.
\end{proof}
\begin{definition}[Mass Quantity]
  \label{def:physics:phyx-mini-0757:phyxminiproblems-problemphyxmini0757-massquantity}
  \lean{PhyXMiniProblems.ProblemPhyXMini0757.MassQuantity}
  A nonnegative, unit-independent physical mass
\end{definition}
\begin{proof}
  This is the declared model definition of Mass Quantity.
\end{proof}
\begin{definition}[Acceleration Quantity]
  \label{def:physics:phyx-mini-0757:phyxminiproblems-problemphyxmini0757-accelerationquantity}
  \lean{PhyXMiniProblems.ProblemPhyXMini0757.AccelerationQuantity}
  \uses{def:physics:phyx-mini-0757:phyxminiproblems-problemphyxmini0757-accelerationdimension}
  A nonnegative, unit-independent acceleration magnitude
\end{definition}
\begin{proof}
  This is the declared model definition of Acceleration Quantity.
\end{proof}
\begin{definition}[Force Quantity]
  \label{def:physics:phyx-mini-0757:phyxminiproblems-problemphyxmini0757-forcequantity}
  \lean{PhyXMiniProblems.ProblemPhyXMini0757.ForceQuantity}
  \uses{def:physics:phyx-mini-0757:phyxminiproblems-problemphyxmini0757-forcedimension}
  A nonnegative, unit-independent force magnitude
\end{definition}
\begin{proof}
  This is the declared model definition of Force Quantity.
\end{proof}
\begin{definition}[mass Readout]
  \label{def:physics:phyx-mini-0757:phyxminiproblems-problemphyxmini0757-massreadout}
  \lean{PhyXMiniProblems.ProblemPhyXMini0757.massReadout}
  \uses{def:physics:phyx-mini-0757:phyxminiproblems-problemphyxmini0757-massquantity}
  Read a physical mass in a selected mass unit
\end{definition}
\begin{proof}
  This is the declared model definition of mass Readout.
\end{proof}
\begin{definition}[acceleration Readout]
  \label{def:physics:phyx-mini-0757:phyxminiproblems-problemphyxmini0757-accelerationreadout}
  \lean{PhyXMiniProblems.ProblemPhyXMini0757.accelerationReadout}
  \uses{def:physics:phyx-mini-0757:phyxminiproblems-problemphyxmini0757-accelerationquantity}
  Read acceleration in coherent selected length and time units
\end{definition}
\begin{proof}
  This is the declared model definition of acceleration Readout.
\end{proof}
\begin{definition}[force Readout]
  \label{def:physics:phyx-mini-0757:phyxminiproblems-problemphyxmini0757-forcereadout}
  \lean{PhyXMiniProblems.ProblemPhyXMini0757.forceReadout}
  \uses{def:physics:phyx-mini-0757:phyxminiproblems-problemphyxmini0757-forcequantity}
  Read force in the coherent unit induced by selected base units
\end{definition}
\begin{proof}
  This is the declared model definition of force Readout.
\end{proof}
\begin{definition}[mass In Kilograms]
  \label{def:physics:phyx-mini-0757:phyxminiproblems-problemphyxmini0757-massinkilograms}
  \lean{PhyXMiniProblems.ProblemPhyXMini0757.massInKilograms}
  \uses{def:physics:phyx-mini-0757:phyxminiproblems-problemphyxmini0757-massquantity, def:physics:phyx-mini-0757:phyxminiproblems-problemphyxmini0757-massreadout}
  Kilogram readout of a physical mass
\end{definition}
\begin{proof}
  This is the declared model definition of mass In Kilograms.
\end{proof}
\begin{definition}[acceleration In Meters Per Second Squared]
  \label{def:physics:phyx-mini-0757:phyxminiproblems-problemphyxmini0757-accelerationinmeterspersecondsquared}
  \lean{PhyXMiniProblems.ProblemPhyXMini0757.accelerationInMetersPerSecondSquared}
  \uses{def:physics:phyx-mini-0757:phyxminiproblems-problemphyxmini0757-accelerationquantity, def:physics:phyx-mini-0757:phyxminiproblems-problemphyxmini0757-accelerationreadout}
  Metre-per-second-squared readout of an acceleration magnitude
\end{definition}
\begin{proof}
  This is the declared model definition of acceleration In Meters Per Second Squared.
\end{proof}
\begin{definition}[force In Newtons]
  \label{def:physics:phyx-mini-0757:phyxminiproblems-problemphyxmini0757-forceinnewtons}
  \lean{PhyXMiniProblems.ProblemPhyXMini0757.forceInNewtons}
  \uses{def:physics:phyx-mini-0757:phyxminiproblems-problemphyxmini0757-forcequantity, def:physics:phyx-mini-0757:phyxminiproblems-problemphyxmini0757-forcereadout}
  Newton readout of a physical force magnitude
\end{definition}
\begin{proof}
  This is the declared model definition of force In Newtons.
\end{proof}
\begin{definition}[Disk Label]
  \label{def:physics:phyx-mini-0757:phyxminiproblems-problemphyxmini0757-disklabel}
  \lean{PhyXMiniProblems.ProblemPhyXMini0757.DiskLabel}
  The four disks, named from top to bottom as in the supplied image
\end{definition}
\begin{proof}
  This is the declared model definition of Disk Label.
\end{proof}
\begin{definition}[Short Cord Label]
  \label{def:physics:phyx-mini-0757:phyxminiproblems-problemphyxmini0757-shortcordlabel}
  \lean{PhyXMiniProblems.ProblemPhyXMini0757.ShortCordLabel}
  The three short cords whose tensions are labelled in the image
\end{definition}
\begin{proof}
  This is the declared model definition of Short Cord Label.
\end{proof}
\begin{definition}[Cord Segment]
  \label{def:physics:phyx-mini-0757:phyxminiproblems-problemphyxmini0757-cordsegment}
  \lean{PhyXMiniProblems.ProblemPhyXMini0757.CordSegment}
  All tension-carrying cord segments, including the long top cord
\end{definition}
\begin{proof}
  This is the declared model definition of Cord Segment.
\end{proof}
\begin{definition}[Figure Object]
  \label{def:physics:phyx-mini-0757:phyxminiproblems-problemphyxmini0757-figureobject}
  \lean{PhyXMiniProblems.ProblemPhyXMini0757.FigureObject}
  Physical objects visible in the supplied raster
\end{definition}
\begin{proof}
  This is the declared model definition of Figure Object.
\end{proof}
\begin{definition}[Figure Text Label]
  \label{def:physics:phyx-mini-0757:phyxminiproblems-problemphyxmini0757-figuretextlabel}
  \lean{PhyXMiniProblems.ProblemPhyXMini0757.FigureTextLabel}
  Literal disk and tension labels visible in the supplied raster
\end{definition}
\begin{proof}
  This is the declared model definition of Figure Text Label.
\end{proof}
\begin{definition}[Suspension State]
  \label{def:physics:phyx-mini-0757:phyxminiproblems-problemphyxmini0757-suspensionstate}
  \lean{PhyXMiniProblems.ProblemPhyXMini0757.SuspensionState}
  The motion state required for the force-balance model
\end{definition}
\begin{proof}
  This is the declared model definition of Suspension State.
\end{proof}
\begin{definition}[short Cord Segment]
  \label{def:physics:phyx-mini-0757:phyxminiproblems-problemphyxmini0757-shortcordsegment}
  \lean{PhyXMiniProblems.ProblemPhyXMini0757.shortCordSegment}
  \uses{def:physics:phyx-mini-0757:phyxminiproblems-problemphyxmini0757-shortcordlabel, def:physics:phyx-mini-0757:phyxminiproblems-problemphyxmini0757-cordsegment}
  Regard a labelled short cord as one of the system's cord segments
\end{definition}
\begin{proof}
  This is the declared model definition of short Cord Segment.
\end{proof}
\begin{definition}[expected Upper Disk]
  \label{def:physics:phyx-mini-0757:phyxminiproblems-problemphyxmini0757-expectedupperdisk}
  \lean{PhyXMiniProblems.ProblemPhyXMini0757.expectedUpperDisk}
  \uses{def:physics:phyx-mini-0757:phyxminiproblems-problemphyxmini0757-disklabel, def:physics:phyx-mini-0757:phyxminiproblems-problemphyxmini0757-shortcordlabel}
  Upper endpoint of each short cord, as read from the figure
\end{definition}
\begin{proof}
  This is the declared model definition of expected Upper Disk.
\end{proof}
\begin{definition}[expected Lower Disk]
  \label{def:physics:phyx-mini-0757:phyxminiproblems-problemphyxmini0757-expectedlowerdisk}
  \lean{PhyXMiniProblems.ProblemPhyXMini0757.expectedLowerDisk}
  \uses{def:physics:phyx-mini-0757:phyxminiproblems-problemphyxmini0757-disklabel, def:physics:phyx-mini-0757:phyxminiproblems-problemphyxmini0757-shortcordlabel}
  Lower endpoint of each short cord, as read from the figure
\end{definition}
\begin{proof}
  This is the declared model definition of expected Lower Disk.
\end{proof}
\begin{definition}[upper Cord]
  \label{def:physics:phyx-mini-0757:phyxminiproblems-problemphyxmini0757-uppercord}
  \lean{PhyXMiniProblems.ProblemPhyXMini0757.upperCord}
  \uses{def:physics:phyx-mini-0757:phyxminiproblems-problemphyxmini0757-disklabel, def:physics:phyx-mini-0757:phyxminiproblems-problemphyxmini0757-cordsegment}
  Cord segment immediately above each disk
\end{definition}
\begin{proof}
  This is the declared model definition of upper Cord.
\end{proof}
\begin{definition}[lower Cord]
  \label{def:physics:phyx-mini-0757:phyxminiproblems-problemphyxmini0757-lowercord}
  \lean{PhyXMiniProblems.ProblemPhyXMini0757.lowerCord}
  \uses{def:physics:phyx-mini-0757:phyxminiproblems-problemphyxmini0757-disklabel, def:physics:phyx-mini-0757:phyxminiproblems-problemphyxmini0757-cordsegment}
  Cord segment immediately below a disk, if one exists
\end{definition}
\begin{proof}
  This is the declared model definition of lower Cord.
\end{proof}
\begin{definition}[Supplied Four Disk Figure]
  \label{def:physics:phyx-mini-0757:phyxminiproblems-problemphyxmini0757-suppliedfourdiskfigure}
  \lean{PhyXMiniProblems.ProblemPhyXMini0757.SuppliedFourDiskFigure}
  \uses{def:physics:phyx-mini-0757:phyxminiproblems-problemphyxmini0757-disklabel, def:physics:phyx-mini-0757:phyxminiproblems-problemphyxmini0757-shortcordlabel, def:physics:phyx-mini-0757:phyxminiproblems-problemphyxmini0757-figureobject, def:physics:phyx-mini-0757:phyxminiproblems-problemphyxmini0757-figuretextlabel}
  Typed evidence transcribed from image 757.png. The raster records the vertical order, cord incidence, and symbolic labels, but does not display a numerical mass for disk C
\end{definition}
\begin{proof}
  This is the declared model definition of Supplied Four Disk Figure.
\end{proof}
\begin{definition}[Four Disk Pulley Setup]
  \label{def:physics:phyx-mini-0757:phyxminiproblems-problemphyxmini0757-fourdiskpulleysetup}
  \lean{PhyXMiniProblems.ProblemPhyXMini0757.FourDiskPulleySetup}
  \uses{def:physics:phyx-mini-0757:phyxminiproblems-problemphyxmini0757-massquantity, def:physics:phyx-mini-0757:phyxminiproblems-problemphyxmini0757-accelerationquantity, def:physics:phyx-mini-0757:phyxminiproblems-problemphyxmini0757-forcequantity, def:physics:phyx-mini-0757:phyxminiproblems-problemphyxmini0757-disklabel, def:physics:phyx-mini-0757:phyxminiproblems-problemphyxmini0757-cordsegment, def:physics:phyx-mini-0757:phyxminiproblems-problemphyxmini0757-suspensionstate, def:physics:phyx-mini-0757:phyxminiproblems-problemphyxmini0757-suppliedfourdiskfigure}
  Independent physical quantities in the four-disk system. Disk weights and cord tensions are fields rather than definitions made from the target mass
\end{definition}
\begin{proof}
  This is the declared model definition of Four Disk Pulley Setup.
\end{proof}
\begin{definition}[Matches Written Scenario]
  \label{def:physics:phyx-mini-0757:phyxminiproblems-problemphyxmini0757-matcheswrittenscenario}
  \lean{PhyXMiniProblems.ProblemPhyXMini0757.MatchesWrittenScenario}
  \uses{def:physics:phyx-mini-0757:phyxminiproblems-problemphyxmini0757-fourdiskpulleysetup}
  Qualitative conditions stated or implicit in the suspended arrangement
\end{definition}
\begin{proof}
  This is the declared model definition of Matches Written Scenario.
\end{proof}
\begin{definition}[Matches Supplied Figure]
  \label{def:physics:phyx-mini-0757:phyxminiproblems-problemphyxmini0757-matchessuppliedfigure}
  \lean{PhyXMiniProblems.ProblemPhyXMini0757.MatchesSuppliedFigure}
  \uses{def:physics:phyx-mini-0757:phyxminiproblems-problemphyxmini0757-expectedupperdisk, def:physics:phyx-mini-0757:phyxminiproblems-problemphyxmini0757-expectedlowerdisk, def:physics:phyx-mini-0757:phyxminiproblems-problemphyxmini0757-fourdiskpulleysetup}
  Geometry and symbolic labels read from the primary image
\end{definition}
\begin{proof}
  This is the declared model definition of Matches Supplied Figure.
\end{proof}
\begin{definition}[Matches Problem Force Readouts]
  \label{def:physics:phyx-mini-0757:phyxminiproblems-problemphyxmini0757-matchesproblemforcereadouts}
  \lean{PhyXMiniProblems.ProblemPhyXMini0757.MatchesProblemForceReadouts}
  \uses{def:physics:phyx-mini-0757:phyxminiproblems-problemphyxmini0757-forceinnewtons, def:physics:phyx-mini-0757:phyxminiproblems-problemphyxmini0757-fourdiskpulleysetup}
  Force magnitudes printed in the problem prose. The decimal values are represented exactly: 58.8 = 294/5 and 9.8 = 49/5
\end{definition}
\begin{proof}
  This is the declared model definition of Matches Problem Force Readouts.
\end{proof}
\begin{definition}[Uses Standard Near Earth Gravity]
  \label{def:physics:phyx-mini-0757:phyxminiproblems-problemphyxmini0757-usesstandardnearearthgravity}
  \lean{PhyXMiniProblems.ProblemPhyXMini0757.UsesStandardNearEarthGravity}
  \uses{def:physics:phyx-mini-0757:phyxminiproblems-problemphyxmini0757-accelerationinmeterspersecondsquared, def:physics:phyx-mini-0757:phyxminiproblems-problemphyxmini0757-fourdiskpulleysetup}
  The standard near-Earth gravitational calibration implicit in the recorded textbook answer. It is separated from values explicitly printed in the problem
\end{definition}
\begin{proof}
  This is the declared model definition of Uses Standard Near Earth Gravity.
\end{proof}
\begin{definition}[Has Physical Parameters]
  \label{def:physics:phyx-mini-0757:phyxminiproblems-problemphyxmini0757-hasphysicalparameters}
  \lean{PhyXMiniProblems.ProblemPhyXMini0757.HasPhysicalParameters}
  \uses{def:physics:phyx-mini-0757:phyxminiproblems-problemphyxmini0757-massinkilograms, def:physics:phyx-mini-0757:phyxminiproblems-problemphyxmini0757-accelerationinmeterspersecondsquared, def:physics:phyx-mini-0757:phyxminiproblems-problemphyxmini0757-forceinnewtons, def:physics:phyx-mini-0757:phyxminiproblems-problemphyxmini0757-fourdiskpulleysetup}
  Positivity conditions selecting a nondegenerate physical suspension
\end{definition}
\begin{proof}
  This is the declared model definition of Has Physical Parameters.
\end{proof}
\begin{definition}[lower Tension Readout]
  \label{def:physics:phyx-mini-0757:phyxminiproblems-problemphyxmini0757-lowertensionreadout}
  \lean{PhyXMiniProblems.ProblemPhyXMini0757.lowerTensionReadout}
  \uses{def:physics:phyx-mini-0757:phyxminiproblems-problemphyxmini0757-forcereadout, def:physics:phyx-mini-0757:phyxminiproblems-problemphyxmini0757-disklabel, def:physics:phyx-mini-0757:phyxminiproblems-problemphyxmini0757-lowercord, def:physics:phyx-mini-0757:phyxminiproblems-problemphyxmini0757-fourdiskpulleysetup}
  Read the force carried by the cord below a disk; the bottom disk has no lower cord and therefore contributes zero lower-cord force to its balance equation
\end{definition}
\begin{proof}
  This is the declared model definition of lower Tension Readout.
\end{proof}
\begin{definition}[Satisfies Suspended Disk Statics Laws]
  \label{def:physics:phyx-mini-0757:phyxminiproblems-problemphyxmini0757-satisfiessuspendeddiskstaticslaws}
  \lean{PhyXMiniProblems.ProblemPhyXMini0757.SatisfiesSuspendedDiskStaticsLaws}
  \uses{def:physics:phyx-mini-0757:phyxminiproblems-problemphyxmini0757-massreadout, def:physics:phyx-mini-0757:phyxminiproblems-problemphyxmini0757-accelerationreadout, def:physics:phyx-mini-0757:phyxminiproblems-problemphyxmini0757-forcereadout, def:physics:phyx-mini-0757:phyxminiproblems-problemphyxmini0757-disklabel, def:physics:phyx-mini-0757:phyxminiproblems-problemphyxmini0757-uppercord, def:physics:phyx-mini-0757:phyxminiproblems-problemphyxmini0757-fourdiskpulleysetup, def:physics:phyx-mini-0757:phyxminiproblems-problemphyxmini0757-lowertensionreadout}
  The governing ideal-cord, weight, and static-equilibrium laws. The final mass of disk C is not a field of this structure. Instead, its value must be derived by applying the generic balance equation to disk C together with the independent T₂, T₃, and gravity readouts
\end{definition}
\begin{proof}
  This is the declared model definition of Satisfies Suspended Disk Statics Laws.
\end{proof}
\begin{definition}[Answer Choice]
  \label{def:physics:phyx-mini-0757:phyxminiproblems-problemphyxmini0757-answerchoice}
  \lean{PhyXMiniProblems.ProblemPhyXMini0757.AnswerChoice}
  Labels of the four mass alternatives printed with the problem
\end{definition}
\begin{proof}
  This is the declared model definition of Answer Choice.
\end{proof}
\begin{definition}[displayed Mass In Kilograms]
  \label{def:physics:phyx-mini-0757:phyxminiproblems-problemphyxmini0757-displayedmassinkilograms}
  \lean{PhyXMiniProblems.ProblemPhyXMini0757.displayedMassInKilograms}
  \uses{def:physics:phyx-mini-0757:phyxminiproblems-problemphyxmini0757-answerchoice}
  Mass in kilograms printed beside each answer label
\end{definition}
\begin{proof}
  This is the declared model definition of displayed Mass In Kilograms.
\end{proof}
\begin{definition}[recorded Dataset Answer]
  \label{def:physics:phyx-mini-0757:phyxminiproblems-problemphyxmini0757-recordeddatasetanswer}
  \lean{PhyXMiniProblems.ProblemPhyXMini0757.recordedDatasetAnswer}
  \uses{def:physics:phyx-mini-0757:phyxminiproblems-problemphyxmini0757-answerchoice}
  Dataset answer metadata, deliberately not used as a theorem premise
\end{definition}
\begin{proof}
  This is the declared model definition of recorded Dataset Answer.
\end{proof}
\begin{definition}[Is Unique Matching Displayed Mass]
  \label{def:physics:phyx-mini-0757:phyxminiproblems-problemphyxmini0757-isuniquematchingdisplayedmass}
  \lean{PhyXMiniProblems.ProblemPhyXMini0757.IsUniqueMatchingDisplayedMass}
  \uses{def:physics:phyx-mini-0757:phyxminiproblems-problemphyxmini0757-massinkilograms, def:physics:phyx-mini-0757:phyxminiproblems-problemphyxmini0757-fourdiskpulleysetup, def:physics:phyx-mini-0757:phyxminiproblems-problemphyxmini0757-answerchoice, def:physics:phyx-mini-0757:phyxminiproblems-problemphyxmini0757-displayedmassinkilograms}
  A choice is the unique displayed mass agreeing with disk C
\end{definition}
\begin{proof}
  This is the declared model definition of Is Unique Matching Displayed Mass.
\end{proof}
% --- Archon declaration topology end ---
% --- Archon physics formalization source end ---
